% original PDF p. 21; continuation of Proposition 3.13
There exists a canonical isomorphism of $S$-group schemes
\[
T\isomto\prod_{\underline{\operatorname{Dyn}}(G)/S}\mathbf G_{\mathrm m,\underline{\operatorname{Dyn}}(G)}.
\]
(Recall (cf. Exp.~II, §~1) that the right-hand side is by definition the $S$-functor which associates to $S'\to S$ the group $\Gm(\underline{\operatorname{Dyn}}(G)\times_S S')$, or, equivalently, $\Gm(\underline{\operatorname{Dyn}}(G_{S'}))$.)
\end{proposition}
\begin{proof}[First proof]
For simplicity let us give it in the adjoint case. Consider the composite morphism
\[
T\to\prod_{\mathscr D/S}T_{\mathscr D}\to\prod_{\mathscr D/S}\mathbf G_{\mathrm m,\mathscr D},
\]
where the first morphism is the canonical morphism and the second is $\prod_{\mathscr D/S}\alpha_{\mathscr D}$ (we have put $\mathscr D=\underline{\operatorname{Dyn}}(G)$). To verify that this morphism is an isomorphism, one may suppose $G$ split; now, in this case, it is nothing other than the morphism $T\to(\mathbf G_{\mathrm m,S})^\Delta$ with components $\alpha$, for $\alpha\in\Delta$, and the latter is an isomorphism (Exp.~XXII, 4.3.8).
\end{proof}
\begin{proof}[Second proof]
By 2.8, 3.5 and 3.11, one may suppose that $G=\operatorname{Q\text{-}\acute Ep}_{S'/S}(\mathscr R)$, $T$ being the canonical maximal torus. Over $S'$, one has by Exp.~XXII 4.3.8 an isomorphism $T_{S'}\isomto(\mathbf G_{\mathrm m,S'})^\Delta$, defined by the simple roots (resp. the simple coroots). The group $E$ acts on the right-hand side by permuting $\Delta$. Now the bundle associated to $S'/S$ by $E\to\Aut(\Delta)$ is $\underline{\operatorname{Dyn}}(G)$ (3.7), and one concludes at once.
\end{proof}
Using Prop.~8.4 of the appendix, one derives:
\begin{corollary}{3.14}\label{XXIV.3.14}
Under the preceding conditions, one has
\[
\mathrm H^1(S,T)\isomto\mathrm H^1(\underline{\operatorname{Dyn}}(G),\Gm)\isomto\Pic(\underline{\operatorname{Dyn}}(G)).
\]
In particular, $\mathrm H^1(S,T)=0$ when $S$ is semi-local.
\end{corollary}
\begin{remark}{3.15}\label{XXIV.3.15}
Let $S$ be a scheme, $G$ (resp. $G'$) a reductive $S$-group, $B\supset T$ (resp. $B'\supset T'$) a Killing pair of $G$ (resp. $G'$), and $u\in\underline{\operatorname{Isomext}}(G,G')(S)$. Put (cf. 2.10)
\[
P=\underline{\operatorname{Isomint}}_u(G,B,T;G',B',T');
\]
it is a principal homogeneous bundle under $T^{\mathrm{ad}}$ (by $(f,t)\mto f\circ\mathrm{int}(t)$).

Let, on the other hand, $\mathscr D=\underline{\operatorname{Dyn}}(G)=\underline{\operatorname{Dyn}}(G')$ (identified by means of $u$ (3.5)), and let $L=\underline{\Isom}_{\mathscr D}(\mathfrak g^{\mathscr D},\mathfrak g'^{\mathscr D})$ be the principal homogeneous bundle under $\mathbf G_{\mathrm m,\mathscr D}$ defined by the invertible $\OO_{\mathscr D}$-module
\[
\underline{\Hom}_{\OO_{\mathscr D}}(\mathfrak g^{\mathscr D},\mathfrak g'^{\mathscr D})=\mathfrak g'^{\mathscr D}\otimes(\mathfrak g^{\mathscr D})^\vee.
\]
Each $f\in P(S')$ defines an isomorphism of $\mathfrak g^{\mathscr D}\otimes_{\OS}\OO_{S'}$ onto $\mathfrak g'^{\mathscr D}\otimes_{\OS}\OO_{S'}$, whence a canonical morphism
\[
P\to\prod_{\mathscr D/S}L.
\]
This morphism is an \emph{isomorphism}, compatible with the isomorphism of acting groups
\[
T^{\mathrm{ad}}\isomto\prod_{\mathscr D/S}\mathbf G_{\mathrm m,\mathscr D}
\]
defined above.
% original PDF p. 22
Indeed, it suffices to verify it in the case where the two groups are pinned, where it is easy.

It follows in particular that in the isomorphism
\[
\mathrm H^1(S,T^{\mathrm{ad}})\isomto\Pic(\underline{\operatorname{Dyn}}(G))
\]
of 3.14, the class of the bundle $P$ is taken to $\operatorname{cl}(\mathfrak g'^{\mathscr D})-\operatorname{cl}(\mathfrak g^{\mathscr D})$. The bundle $P$ is therefore trivial if and only if $\mathfrak g'^{\mathscr D}$ and $\mathfrak g^{\mathscr D}$ are isomorphic.

If $(G,B,T)$ is quasi-splittable, for example if for $G$ one takes the group $G'_{\text{q-pin.}}$, with its canonical Killing pair, it follows that the image of the class of $P$ is nothing other than the obstruction to the quasi-splitting of $G'$ defined in 3.9.
\end{remark}
\numpara{3.16}\label{XXIV.3.16}\textbf{Symmetries}
\numpara{3.16.1}\label{XXIV.3.16.1}
Let $S$ be a scheme, $G$ a reductive $S$-group, and $B\supset T$ a Killing pair of $G$. Put $\mathscr D=\underline{\operatorname{Dyn}}(G)$. Recall (Exp.~XXII, 5.9.1) that there exists a unique Borel subgroup $B^-$ of $G$ such that $B\cap B^-=T$. If $X\in\Gamma(\mathscr D,\mathfrak g^{\mathscr D})^\times$ defines a quasi-pinning of $G$ relative to $(B,T)$ (3.9), then $Y=-X^{-1}\in\Gamma(\mathscr D,(\mathfrak g^{\mathscr D})^\vee)^\times$ defines a quasi-pinning of $G$ relative to $(B^-,T)$; one says that it is the \emph{opposite quasi-pinning}.

If $\mathscr R$ is a reduced pinned root datum and $\mathscr E$ an $\mathscr R$-pinning of the reductive $S$-group $G$, one defines an $\mathscr R$-pinning $\mathscr E^-$ said to be \emph{opposite to $\mathscr E$} as follows: one keeps the same maximal torus $T$, takes the opposite of the isomorphism $\mathrm D_S(M)\isomto T$, and ``pins'' by $Y_\alpha=-X_\alpha^{-1}\in\Gamma(S,\mathfrak g^{-\alpha})^\times$, for $\alpha\in\Delta$. The quasi-pinning underlying $\mathscr E^-$ is the quasi-pinning opposite to the quasi-pinning underlying $\mathscr E$.
\begin{remark*}
In the notation of Exp.~XIX, 3.1, if one puts
\[
w_\alpha(X_\alpha)=\exp(X_\alpha)\exp(-X_\alpha^{-1})\exp(X_\alpha),
\]
one has $w_\alpha(X_\alpha)=w_{-\alpha}(Y_\alpha)$ (loc. cit., 3.1 (vi)).
\end{remark*}
\begin{proposition}{3.16.2}\label{XXIV.3.16.2}
Let $S$ be a scheme and $G$ a reductive $S$-group.
\begin{enumeratei}
\item Let $T$ be a maximal torus of $G$; there exists a unique
\[
i_T\in\bigl(\underline{\Aut}_{S\text{-gr.}}(G,T)/T^{\mathrm{ad}}\bigr)(S)\subset\Aut_{S\text{-gr.}}(T)
\]
such that $i_T(t)=t^{-1}$ for every section $t$ of $T$.
\item Let $(B,T)$ be a Killing pair of $G$; there exists a unique section
\[
w_{B,T}\in(\uNorm_G(T)/T)(S)=\mathrm W_G(T)(S)
\]
such that $\mathrm{int}(w_{B,T})(B)=B^-$ (with the evident abuse of language).
\item Let $\mathscr Q=(B,T,X)$ be a quasi-pinning of $G$, and $\mathscr Q^-=(B^-,T,Y)$ the opposite quasi-pinning; there exists a unique inner automorphism of $G$
\[
n_{\mathscr Q}\in\uNorm_{\ad(G)}(T^{\mathrm{ad}})(S)\subset\Aut_{S\text{-gr.}}(G)
\]
such that $n_{\mathscr Q}(\mathscr Q)=\mathscr Q^-$, i.e. $n_{\mathscr Q}(T)=T$, $n_{\mathscr Q}(B)=B^-$, $n_{\mathscr Q}(X)=Y$.
% original PDF p. 23
\item Let $(\mathscr R,\mathscr E)$ be a pinning of $G$, and $(\mathscr R,\mathscr E^-)$ the opposite pinning; there exists a unique automorphism of $G$
\[
u_{\mathscr E}\in\Aut_{S\text{-gr.}}(G,T)\subset\Aut_{S\text{-gr.}}(G)
\]
such that $u_{\mathscr E}(\mathscr E)=\mathscr E^-$, i.e. $u_{\mathscr E}(t)=t^{-1}$ for every section $t$ of $T$, and $\Ad(u_{\mathscr E})X_\alpha=Y_\alpha$ for every $\alpha\in\Delta$.
\end{enumeratei}
\end{proposition}
\begin{proof}
(ii) follows from Exp.~XXII, 5.5.5 (ii), then (iii) follows from 3.10, and (iv) follows from Exp.~XXIII, 4.1. Finally, to prove (i), one may suppose $G$ pinned. Existence follows from (iv), for example, and uniqueness from the fact that an automorphism of $G$ inducing the identity on $T$ is given by a section of $T^{\mathrm{ad}}$ (2.11).
\end{proof}
\begin{corollary}{3.16.3}\label{XXIV.3.16.3}
One has
\[
i_T^2=w_{B,T}^2=n_{\mathscr Q}^2=u_{\mathscr E}^2=e.
\]
Moreover, $i_T$ (resp. $u_{\mathscr E}$) is $\ne e$ if $G\ne e$, and $w_{B,T}$ (resp. $n_{\mathscr Q}$) is $\ne e$ if $G$ is not a torus.
\end{corollary}
\begin{corollary}{3.16.4}\label{XXIV.3.16.4}
In the situation of (iii) (resp. (iv)), $n_{\mathscr Q}$ projects to $w_{B,T}$ (resp. $u_{\mathscr E}$ projects to $i_T$) under the canonical morphism
\[
\uNorm_{\ad(G)}(T^{\mathrm{ad}})\to\mathrm W_{\ad(G)}(T^{\mathrm{ad}})\simeq\mathrm W_G(T),
\]
resp.
\[
\underline{\Aut}_{S\text{-gr.}}(G,T)\to\underline{\Aut}_{S\text{-gr.}}(G,T)/T^{\mathrm{ad}}.
\]
\end{corollary}
\begin{corollary}{3.16.5}\label{XXIV.3.16.5}
The preceding definitions are compatible with extension of the base, and are functorial under isomorphism (in an evident sense).
\end{corollary}
\begin{proposition}{3.16.6}\label{XXIV.3.16.6}
\begin{enumeratei}
\item One can define uniquely for every reductive group $G$ over a scheme $S$ an element
\[
s_G\in\underline{\operatorname{Autext}}(G)(S)
\]
in such a way that this construction is functorial in $G$ under isomorphism, is compatible with base change, and whenever $T$ is a maximal torus of the reductive $S$-group $G$, $s_G$ is the image of $i_T$ under the canonical morphism
\[
\underline{\Aut}_{S\text{-gr.}}(G,T)/T^{\mathrm{ad}}\to\underline{\operatorname{Autext}}(G).
\]
\item One has $s_G^2=e$, and $s_G$ is a central element of $\underline{\operatorname{Autext}}(G)$.
\item Under the conditions of 3.16.2 (ii), if one identifies $\underline{\Aut}_{S\text{-gr.}}(G,B,T)/T^{\mathrm{ad}}$ with $\underline{\operatorname{Autext}}(G)$ (2.2), one has
\[
w_{B,T}i_T=i_Tw_{B,T}=s_G.
\]
\item Under the conditions of 3.16.2 (iv), if one identifies $\Aut_{S\text{-gr.-pin}}(G,\mathscr R,\mathscr E)$ with $\underline{\operatorname{Autext}}(G)$ (1.3 (iii)), one has
\[
n_{\mathscr E}u_{\mathscr E}=u_{\mathscr E}n_{\mathscr E}=s_G.
\]
\end{enumeratei}
\end{proposition}
% original PDF p. 24
\begin{proof}
(i) is proved without difficulty by descent. Moreover, since $i_T$ is evidently a central section of square $e$ in $\Aut_{S\text{-gr.}}(T)$, (ii) follows immediately; (iii) is a consequence of (iv) by descent. Finally, under the conditions of (iv), it is clear that $n_{\mathscr E}u_{\mathscr E}=u_{\mathscr E}n_{\mathscr E}$ and that this automorphism of $G$ respects the pinning; under the identification made, it is therefore equal to its image in $\underline{\operatorname{Autext}}(G)$; but $n_{\mathscr E}$ is inner and $u_{\mathscr E}$ projects to $s_G$.
\end{proof}
\begin{remark}{3.16.7}\label{XXIV.3.16.7}
(i) One determines $s_G$ explicitly in each case of the classification by means of (iii): for each irreducible pinned root datum, it suffices to compose the symmetry about the origin with the symmetry in the Weyl group (i.e. the element $w_0$ of the Weyl group such that $w_0(\Delta)=-\Delta$). One finds the following results: $s_G=e$ except for $A_n$ ($n\geq2$), $D_n$ ($n$ odd) and $E_6$, in which case $s_G$ is the unique nontrivial ``outer automorphism''.

(ii) The automorphism $u_{\mathscr E}$ is the one used to construct ``the compact real forms'' in the theory of semisimple Lie algebras.
\end{remark}
\begin{remark}{3.16.8}\label{XXIV.3.16.8}
In 3.16.1 one has defined an involution on the $S$-scheme
\[
Q=\underline{\operatorname{Isomint}}(G_{\text{q-pin.}},G)
\]
of quasi-pinnings of $G$ (cf. 3.12.1); by transport of structure from $G_{\text{q-pin.}}$ to $G$, one sees at once that this involution is given by the action of an element of $\ad(G_{\text{q-pin.}})(S)$: the element $n_0$ defined (3.16.2 (iii)) by the canonical quasi-pinning of $G_{\text{q-pin.}}$.

In the same way, one has defined an involution on the $S$-scheme
\[
P=\underline{\Isom}_{S\text{-gr.}}(\operatorname{\acute Ep}_S(\mathscr R),G)
\]
of $\mathscr R$-pinnings of $G$ (cf. 1.20). Reasoning as before, one sees that this involution is given by the action of the automorphism $u_0$ of $\operatorname{\acute Ep}_S(\mathscr R)$ defined (3.16.2 (iv)) by the canonical pinning of $\operatorname{\acute Ep}_S(\mathscr R)$.
\end{remark}

\sgasection{4}{Isotriviality of reductive groups and principal bundles under reductive groups}
\numpara{4.1}\label{XXIV.4.1}\textbf{Definitions. Isotriviality theorem}
\begin{definition}{4.1.1}\label{XXIV.4.1.1}
Let $S$ be a scheme, $G$ an $S$-group scheme, and $P$ a principal homogeneous bundle under $G$. One says that $P$ is \emph{locally isotrivial} (resp. \emph{semi-locally isotrivial}) if for every point $s\in S$ (resp. every finite set $F$ of points of $S$ contained in an affine open subset) there exist an open subset $U$ of $S$ containing $s$ (resp. $F$) and a finite surjective étale morphism $S'\to U$ such that $P_{S'}$ is trivial.
\end{definition}
\begin{definition}{4.1.2}\label{XXIV.4.1.2}
Let $S$ be a scheme and $G$ a reductive $S$-group. One says that $G$ is locally isotrivial (resp. semi-locally isotrivial) if for every point $s\in S$ (resp. every finite set $F$ of points of $S$ contained in an affine open subset) there exist an open subset $U$ of $S$ containing $s$ (resp. $F$) and a finite surjective étale morphism $S'\to U$ such that $G_{S'}$ is splittable.
\end{definition}
\begin{remark}{4.1.3}\label{XXIV.4.1.3}
(i) The equivalence of categories of 1.17 respects, by definition, local (resp. semi-local) isotriviality.
% original PDF p. 25

(ii) Add to the conditions of 4.1.1: $G$ locally of finite presentation over $S$. Then the principal bundle $P$ (or the reductive group $G$) is locally isotrivial (resp. semi-locally isotrivial) if and only if for every $S'\to S$, with $S'$ local (resp. semi-local), $P_{S'}$ is isotrivial (or $G_{S'}$ is isotrivial), i.e. if there exists a finite surjective étale morphism $S''\to S'$ such that $P_{S''}$ is trivial (or $G_{S''}$ split).

(iii) In the case of tori, Definition 4.1.2 coincides with that of Exp.~IX, 1.1.
\end{remark}
\numpara{4.1.4}\label{XXIV.4.1.4}
Recall (Exp.~XXII, 4.3 and 6.2) that for every reductive group $G$ we have introduced the groups $\operatorname{rad}(G)$, $\operatorname{corad}(G)$ and $\operatorname{der}(G)$. The groups $\operatorname{rad}(G)$ and $\operatorname{corad}(G)$ are tori, which are split when $G$ is split; moreover, there exists an isogeny $\operatorname{rad}(G)\to\operatorname{corad}(G)$. The $S$-group $\operatorname{der}(G)$ is semisimple, and one has $G/\operatorname{der}(G)=\operatorname{corad}(G)$; it follows that for every principal homogeneous bundle $P$ under $G$, $P/\operatorname{der}(G)$ is a principal homogeneous bundle under $\operatorname{corad}(G)$. This being said, one has:
\begin{theorem}{4.1.5}\label{XXIV.4.1.5}
Let $S$ be a scheme and $G$ a reductive $S$-group of constant type.
\begin{enumeratei}
\item The following conditions are equivalent:
\begin{enumerate}[label=\textup{(\alph*)}]
\item $G$ is locally (resp. semi-locally) isotrivial.
\item The torus $\operatorname{rad}(G)$ is so.
\item[(b')] The torus $\operatorname{corad}(G)$ is so.
\item The Galois covering of $S$ associated to $G$ (3.11) is so.
\end{enumerate}
If $T$ is a maximal torus of $G$, these conditions are also equivalent to
\begin{enumerate}[label=\textup{(\alph*)},start=4]
\item The torus $T$ is locally (resp. semi-locally) isotrivial.
\end{enumerate}
\item Let $P$ be a principal homogeneous bundle under $G$; for $P$ to be locally (resp. semi-locally) isotrivial, it is necessary and sufficient that the principal $\operatorname{corad}(G)$-bundle $P/\operatorname{der}(G)$ be so.
\end{enumeratei}
\end{theorem}
\begin{corollary}{4.1.6}\label{XXIV.4.1.6}
The conditions of (i) hold when $G$ is semisimple or when $S$ is locally noetherian and normal (or, more generally, geometrically unibranch). The conditions of (ii) hold when $G$ is locally (resp. semi-locally) isotrivial.
\end{corollary}
\begin{proof}
For (i), the first assertion is trivial from (b), and the second follows from (c) and Exp.~X, 5.14 and 5.15. For (ii), it suffices to observe that, by theorem 90, a principal bundle under a split torus is semi-locally isotrivial.
\end{proof}
\numpara{4.2}\label{XXIV.4.2}\textbf{Proof: the semisimple case}
Let us first prove, for later reference:
\begin{proposition}{4.2.1}\label{XXIV.4.2.1}
Let $S$ be a scheme, $G$ a reductive $S$-group, $T$ a maximal torus of $G$ (resp. $B$ a Borel subgroup, resp. $B\supset T$ a Killing pair of $G$), $P$ a principal homogeneous bundle under $G$, and $G'$ the twisted form of $G$ associated to $P$ (via the morphism $\mathrm{int}:G\to\underline{\Aut}_{S\text{-gr.}}(G)$). One has canonical isomorphisms
\[
P/\uNorm_G(T)\isomto\underline{\operatorname{Tor}}(G'),\qquad
P/B\isomto\underline{\operatorname{Bor}}(G'),\qquad
P/T\isomto\underline{\operatorname{Kil}}(G').
\]
\end{proposition}
% original PDF p. 26
\begin{proof}
By construction, $G'$ is the quotient of $P\times_S G$ by a certain action of $G$ ($(p,g')g=(pg,g^{-1}g'g)$); one therefore has a morphism $P\times_S G\to G'$, i.e. a morphism
\[
P\to\underline{\Hom}_S(G,G'),
\]
which, as one sees at once, factors through a morphism
\[
f:P\to\underline{\Isom}_{S\text{-gr.}}(G,G'),
\]
(to verify this assertion, one may suppose $G=P$, in which case one has $G'=G$, $f=\mathrm{int}$). The set-theoretic map $p\mto f(p)(T)$ defines a morphism
\[
P\to\underline{\operatorname{Tor}}(G').
\]
To verify that this morphism induces an isomorphism $P/\uNorm_G(T)\isomto\underline{\operatorname{Tor}}(G')$ as asserted, one may again suppose $P=G$, in which case one is reduced to Exp.~XXII, 5.8.3 (iii). (In fact, loc. cit. ought to be replaced by the above statement.) One reasons likewise for $\underline{\operatorname{Bor}}$ and $\underline{\operatorname{Kil}}$.
\end{proof}
\begin{proposition}{4.2.2}\label{XXIV.4.2.2}
Let $S$ be a semi-local scheme and $G$ a semisimple $S$-group of constant type.
\begin{enumeratei}
\item $G$ is isotrivial.
\item Every principal bundle under $G$ is isotrivial.
\end{enumeratei}
\end{proposition}
Let us prove (i). After a finite surjective étale extension of the base, one may, by 3.9.2, suppose $G$ quasi-split. But then $G$ is isotrivial by construction (3.11, the group $E$ is finite). To prove (ii), one may, by (i), suppose $G$ split; one is then reduced to:
\begin{lemma}{4.2.3}\label{XXIV.4.2.3}
Let $S$ be a semi-local scheme. Every principal bundle under a split reductive group is isotrivial.
\end{lemma}
\begin{proof}
Indeed, with the notation of 4.2.1, where $B\supset T$ denotes the canonical Killing pair of the split $G$, the $S$-scheme $\underline{\operatorname{Kil}}(G')$ has a section, after a finite surjective étale extension of the base, by 3.9.2. One can therefore reduce the structure group of $G$ to $T$; now $T$ is split, hence $\mathrm H^1(S,T)=0$ (theorem 90).
% typo?: source says structure group of G rather than P; kept as printed.
\end{proof}
\begin{corollary}{4.2.4}\label{XXIV.4.2.4}
Let $S$ be a scheme and
\[
1\to G\to G'\to G''\to1
\]
an exact sequence of $S$-group schemes (for (fpqc)), $G$ being semisimple of constant type. Let $P$ be a principal homogeneous bundle under $G'$, and suppose the associated sheaf\nde{We have replaced ``bundle'' by ``sheaf'', and then $G'$ by $G''$.} $P/G$ representable (for example $G''$ affine over $S$). For $P$ to be locally isotrivial (resp. semi-locally isotrivial), it is necessary and sufficient that $P/G$ be so (as a bundle under $G''$\footnotemark[24]).
\end{corollary}
\begin{proof}
If $P$ is trivial, $P/G$ is also trivial, which shows that the condition is necessary. Conversely, suppose $S$ local (resp. semi-local) and $P/G$ isotrivial, hence trivialized by an étale finite surjective extension $S'$ of $S$. Extending the base to $S'$, one can reduce the structure group of $P_{S'}$ to $G_{S'}$.
% Proof continues in en-05.tex; last sentence finishes on p. 27.
